% Options for packages loaded elsewhere
\PassOptionsToPackage{unicode}{hyperref}
\PassOptionsToPackage{hyphens}{url}
\documentclass[
]{article}
\usepackage{xcolor}
\usepackage[margin=1in]{geometry}
\usepackage{amsmath,amssymb}
\setcounter{secnumdepth}{5}
\usepackage{iftex}
\ifPDFTeX
  \usepackage[T1]{fontenc}
  \usepackage[utf8]{inputenc}
  \usepackage{textcomp} % provide euro and other symbols
\else % if luatex or xetex
  \usepackage{unicode-math} % this also loads fontspec
  \defaultfontfeatures{Scale=MatchLowercase}
  \defaultfontfeatures[\rmfamily]{Ligatures=TeX,Scale=1}
\fi
\usepackage{lmodern}
\ifPDFTeX\else
  % xetex/luatex font selection
\fi
% Use upquote if available, for straight quotes in verbatim environments
\IfFileExists{upquote.sty}{\usepackage{upquote}}{}
\IfFileExists{microtype.sty}{% use microtype if available
  \usepackage[]{microtype}
  \UseMicrotypeSet[protrusion]{basicmath} % disable protrusion for tt fonts
}{}
\makeatletter
\@ifundefined{KOMAClassName}{% if non-KOMA class
  \IfFileExists{parskip.sty}{%
    \usepackage{parskip}
  }{% else
    \setlength{\parindent}{0pt}
    \setlength{\parskip}{6pt plus 2pt minus 1pt}}
}{% if KOMA class
  \KOMAoptions{parskip=half}}
\makeatother
\usepackage{longtable,booktabs,array}
\usepackage{calc} % for calculating minipage widths
% Correct order of tables after \paragraph or \subparagraph
\usepackage{etoolbox}
\makeatletter
\patchcmd\longtable{\par}{\if@noskipsec\mbox{}\fi\par}{}{}
\makeatother
% Allow footnotes in longtable head/foot
\IfFileExists{footnotehyper.sty}{\usepackage{footnotehyper}}{\usepackage{footnote}}
\makesavenoteenv{longtable}
\usepackage{graphicx}
\makeatletter
\newsavebox\pandoc@box
\newcommand*\pandocbounded[1]{% scales image to fit in text height/width
  \sbox\pandoc@box{#1}%
  \Gscale@div\@tempa{\textheight}{\dimexpr\ht\pandoc@box+\dp\pandoc@box\relax}%
  \Gscale@div\@tempb{\linewidth}{\wd\pandoc@box}%
  \ifdim\@tempb\p@<\@tempa\p@\let\@tempa\@tempb\fi% select the smaller of both
  \ifdim\@tempa\p@<\p@\scalebox{\@tempa}{\usebox\pandoc@box}%
  \else\usebox{\pandoc@box}%
  \fi%
}
% Set default figure placement to htbp
\def\fps@figure{htbp}
\makeatother
\setlength{\emergencystretch}{3em} % prevent overfull lines
\providecommand{\tightlist}{%
  \setlength{\itemsep}{0pt}\setlength{\parskip}{0pt}}
\usepackage{bookmark}
\IfFileExists{xurl.sty}{\usepackage{xurl}}{} % add URL line breaks if available
\urlstyle{same}
\hypersetup{
  pdftitle={Rapport d'analyse agroSAR},
  pdfauthor={EL HENDAOUI Aya},
  hidelinks,
  pdfcreator={LaTeX via pandoc}}

\title{Rapport d'analyse agroSAR}
\usepackage{etoolbox}
\makeatletter
\providecommand{\subtitle}[1]{% add subtitle to \maketitle
  \apptocmd{\@title}{\par {\large #1 \par}}{}{}
}
\makeatother
\subtitle{Suivi agricole par radar Sentinel-1 et SoilGrids}
\author{EL HENDAOUI Aya}
\date{14 juin 2026}

\begin{document}
\maketitle

{
\setcounter{tocdepth}{2}
\tableofcontents
}
\section{Introduction}\label{introduction}

Ce rapport présente les résultats de l'analyse des données radar
\textbf{Sentinel-1} pour le suivi agricole, couplées aux données
d'humidité du sol issues de \textbf{SoilGrids (ISRIC)}.

\begin{center}\rule{0.5\linewidth}{0.5pt}\end{center}

\section{Données utilisées}\label{donnuxe9es-utilisuxe9es}

\begin{longtable}[]{@{}llll@{}}
\toprule\noalign{}
Source & Produit & Résolution & Accès \\
\midrule\noalign{}
\endhead
\bottomrule\noalign{}
\endlastfoot
Copernicus / ESA & Sentinel-1 GRD VV/VH & 10 m & Gratuit \\
ISRIC SoilGrids & Humidité volumétrique 0--5 cm & 250 m & Gratuit \\
\end{longtable}

\begin{center}\rule{0.5\linewidth}{0.5pt}\end{center}

\section{Analyse du signal SAR}\label{analyse-du-signal-sar}

\subsection{Évolution temporelle du
backscatter}\label{uxe9volution-temporelle-du-backscatter}

\begin{figure}

{\centering \includegraphics{C:\Users\EL HENDAOUI AYA\Desktop\agroSAR\docs\rapport_agroSAR_files/figure-latex/timeseries-1} 

}

\caption{Évolution temporelle des bandes VV et VH}\label{fig:timeseries}
\end{figure}

\subsection{Distribution du backscatter VV et
VH}\label{distribution-du-backscatter-vv-et-vh}

\begin{figure}

{\centering \includegraphics{C:\Users\EL HENDAOUI AYA\Desktop\agroSAR\docs\rapport_agroSAR_files/figure-latex/distribution-1} 

}

\caption{Distribution des valeurs de backscatter}\label{fig:distribution}
\end{figure}

\begin{center}\rule{0.5\linewidth}{0.5pt}\end{center}

\section{Humidité du sol ---
SoilGrids}\label{humidituxe9-du-sol-soilgrids}

\begin{longtable}[]{@{}rrllrl@{}}
\caption{Données d'humidité du sol --- SoilGrids}\tabularnewline
\toprule\noalign{}
lon & lat & depth & property & value & unit \\
\midrule\noalign{}
\endfirsthead
\toprule\noalign{}
lon & lat & depth & property & value & unit \\
\midrule\noalign{}
\endhead
\bottomrule\noalign{}
\endlastfoot
-5.0 & 33.9 & 0-5cm & wv0033 & 28.4 & \%vol \\
-5.2 & 33.7 & 0-5cm & wv0033 & 31.2 & \%vol \\
-5.4 & 33.5 & 0-5cm & wv0033 & 25.7 & \%vol \\
\end{longtable}

\begin{center}\rule{0.5\linewidth}{0.5pt}\end{center}

\section{Modèle d'estimation de
l'humidité}\label{moduxe8le-destimation-de-lhumidituxe9}

\subsection{Performance du modèle}\label{performance-du-moduxe8le}

\begin{verbatim}
## RMSE : 4.2
\end{verbatim}

\begin{verbatim}
## R²   : 0.87
\end{verbatim}

\subsection{Graphique observé vs
prédit}\label{graphique-observuxe9-vs-pruxe9dit}

\begin{figure}

{\centering \includegraphics{C:\Users\EL HENDAOUI AYA\Desktop\agroSAR\docs\rapport_agroSAR_files/figure-latex/validation-1} 

}

\caption{Validation : valeurs observées vs prédites}\label{fig:validation}
\end{figure}

\subsection{Importance des variables}\label{importance-des-variables}

\begin{figure}

{\centering \includegraphics{C:\Users\EL HENDAOUI AYA\Desktop\agroSAR\docs\rapport_agroSAR_files/figure-latex/importance-1} 

}

\caption{Importance des variables — Random Forest}\label{fig:importance}
\end{figure}

\begin{center}\rule{0.5\linewidth}{0.5pt}\end{center}

\section{Conclusion}\label{conclusion}

Ce rapport montre que les données radar Sentinel-1 permettent :

\begin{itemize}
\tightlist
\item
  Le suivi de l'\textbf{évolution temporelle} des cultures via VV et VH
\item
  L'\textbf{estimation de l'humidité du sol} avec Random Forest
\item
  La \textbf{classification des stades culturaux} à partir des indices
  SAR
\end{itemize}

\begin{center}\rule{0.5\linewidth}{0.5pt}\end{center}

\section{Informations de session}\label{informations-de-session}

\begin{verbatim}
## R version 4.5.1 (2025-06-13 ucrt)
## Platform: x86_64-w64-mingw32/x64
## Running under: Windows 11 x64 (build 26200)
## 
## Matrix products: default
##   LAPACK version 3.12.1
## 
## locale:
## [1] LC_COLLATE=French_France.utf8  LC_CTYPE=French_France.utf8   
## [3] LC_MONETARY=French_France.utf8 LC_NUMERIC=C                  
## [5] LC_TIME=French_France.utf8    
## 
## time zone: Europe/Paris
## tzcode source: internal
## 
## attached base packages:
## [1] stats     graphics  grDevices utils     datasets  methods   base     
## 
## other attached packages:
## [1] agroSAR_0.1.0  dplyr_1.2.1    ggplot2_4.0.3  terra_1.9-27   testthat_3.3.2
## 
## loaded via a namespace (and not attached):
##   [1] RColorBrewer_1.1-3   sys_3.4.3            rstudioapi_0.18.0   
##   [4] jsonlite_2.0.0       wk_0.9.5             magrittr_2.0.5      
##   [7] farver_2.1.2         rmarkdown_2.31       fs_2.1.0            
##  [10] ragg_1.5.2           vctrs_0.7.3          memoise_2.0.1       
##  [13] askpass_1.2.1        base64enc_0.1-6      gh_1.5.0            
##  [16] htmltools_0.5.9      leafsync_0.1.0       usethis_3.2.1       
##  [19] curl_7.1.0           raster_3.6-32        s2_1.1.9            
##  [22] sass_0.4.10          bslib_0.11.0         KernSmooth_2.23-26  
##  [25] htmlwidgets_1.6.4    desc_1.4.3           httr2_1.2.2         
##  [28] plotly_4.12.0        stars_0.7-2          cachem_1.1.0        
##  [31] commonmark_2.0.0     whisker_0.4.1        lifecycle_1.0.5     
##  [34] pkgconfig_2.0.3      Matrix_1.7-3         cols4all_0.10       
##  [37] R6_2.6.1             fastmap_1.2.0        rcmdcheck_1.4.0     
##  [40] digest_0.6.39        colorspace_2.1-2     rprojroot_2.1.1     
##  [43] leafem_0.2.5         pkgload_1.5.2        textshaping_1.0.5   
##  [46] crosstalk_1.2.2      lwgeom_0.2-16        labeling_0.4.3      
##  [49] randomForest_4.7-1.2 spacesXYZ_1.6-0      httr_1.4.8          
##  [52] abind_1.4-8          mgcv_1.9-3           compiler_4.5.1      
##  [55] microbenchmark_1.5.0 proxy_0.4-29         withr_3.0.2         
##  [58] S7_0.2.2             DBI_1.3.0            logger_0.4.2        
##  [61] pak_0.9.5            pkgbuild_1.4.8       maptiles_0.11.0     
##  [64] openssl_2.4.1        rappdirs_0.3.4       sessioninfo_1.2.3   
##  [67] tmaptools_3.3        leaflet_2.2.3        classInt_0.4-11     
##  [70] tools_4.5.1          units_1.0-1          otel_0.2.0          
##  [73] tmap_4.3             leaflegend_1.2.8     clipr_0.8.0         
##  [76] glue_1.8.1           callr_3.7.6          nlme_3.1-168        
##  [79] grid_4.5.1           sf_1.1-1             generics_0.1.4      
##  [82] gtable_0.3.6         class_7.3-23         tidyr_1.3.2         
##  [85] data.table_1.18.2.1  sp_2.2-1             xml2_1.5.2          
##  [88] pillar_1.11.1        splines_4.5.1        lattice_0.22-7      
##  [91] tidyselect_1.2.1     rstac_1.0.1          knitr_1.51          
##  [94] gitcreds_0.1.2       xfun_0.57            devtools_2.5.2      
##  [97] credentials_2.0.3    brio_1.1.5           lazyeval_0.2.3      
## [100] xopen_1.0.1          yaml_2.3.12          evaluate_1.0.5      
## [103] codetools_0.2-20     tibble_3.3.1         cli_3.6.6           
## [106] systemfonts_1.3.2    jquerylib_0.1.4      processx_3.9.0      
## [109] roxygen2_8.0.0       Rcpp_1.1.1-1.1       gert_2.3.1          
## [112] png_0.1-9            XML_3.99-0.23        parallel_4.5.1      
## [115] ellipsis_0.3.3       prettyunits_1.2.0    jpeg_0.1-11         
## [118] viridisLite_0.4.3    scales_1.4.0         e1071_1.7-17        
## [121] purrr_1.2.2          crayon_1.5.3         rlang_1.2.0         
## [124] waldo_0.6.2
\end{verbatim}

\end{document}
